\documentclass[11pt]{article}

\usepackage{amsmath,amssymb}
\usepackage{xurl}
\usepackage[hidelinks]{hyperref}
\setlength{\emergencystretch}{2em}

\input{command}
\renewcommand{\ghissue}[1]{%
  \href{https://github.com/Dengnifer/MIPStarRE-A/issues/#1}{GitHub issue \##1}}

\title{The Decoding Identity in the Pauli Extraction Argument}
\author{}
\date{2026-09-03}

\begin{document}
\maketitle

\section{At a glance}

\begin{itemize}
  \item \textbf{Difficulty.} Medium.  The source applies a decoder--encoder
  identity to every bounded individual-degree polynomial, although the identity
  is guaranteed for encoded polynomials and can fail outside the encoding image.
  \item \textbf{Estimated weight.} One restricted algebraic identity and one
  probabilistic estimate showing that the constructed measurement assigns
  little state-dependent mass to outcomes outside the encoding image.
  \item \textbf{Mathlib/project split.} Evaluation of the multilinear
  interpolant is elementary polynomial algebra.  The non-encoding-mass bound
  uses project-specific consistency estimates and Schwartz--Zippel.
  \item \textbf{Key mathematical inputs.} Multilinear interpolation on the
  Boolean cube, injectivity of the low-degree encoding, measurement
  consistency, and Schwartz--Zippel.  The correction is tracked in
  \ghissue{19}.
\end{itemize}

\section{Key theorem forms}

For a word $h\colon\{0,1\}^m\to\mathbb F_q$, let $g_h$ be its multilinear
extension.  Define the chapter's decoder by
\[
  \operatorname{Dec}(g)(y)=g(y),
  \qquad y\in\{0,1\}^m.
\]
Then
\[
  \operatorname{Dec}(g_h)=h
  \quad\text{and}\quad
  \operatorname{Dec}(g_h)\mathbin{\cdot}\operatorname{ind}_m(u)=g_h(u).
\]
More generally, the second identity holds for a polynomial outcome $g$ under
the explicit hypothesis that $g$ belongs to the image of the encoding; no
converse is asserted.

For an arbitrary bounded individual-degree polynomial $g$, the Boolean values
determine its multilinear interpolant, not necessarily $g$ itself.  The
unrestricted identity
\[
  \operatorname{Dec}(g)\mathbin{\cdot}\operatorname{ind}_m(u)=g(u)
\]
is therefore false.

\section{Counterexample and source comparison}

The low-degree encoding is defined at lines~873--897 of
\path{references/qpbt-paper/04_preliminaries.tex}.  The general decoder at
lines~917--924 retains evaluations lying in a chosen subset $H$; the extraction
chapter instead specifies at lines~1419--1420 of
\path{references/qpbt-paper/14_analysis_of_the_pauli_basis_test.tex} that the
decoder returns all evaluations on the Boolean cube.  This is the specialization
$H=\mathbb F_q$, rather than the Boolean-valued shorthand used earlier
\cite[lines~873--924 and 1419--1420]{Ji2020MIPStar}.

Take $m=1$, $d\geq2$, and a field of characteristic two with more than two
elements.  For $g(x)=x^2$, the Boolean evaluations are $g(0)=0$ and $g(1)=1$.
Consequently the decoder is the word whose multilinear extension is $x$.
For any $u\notin\{0,1\}$ with $u^2\neq u$,
\[
  \operatorname{Dec}(g)\mathbin{\cdot}\operatorname{ind}_1(u)=u
  \neq u^2=g(u).
\]
Thus the equality asserted for every $g$ at lines~1483--1492 of the chapter-14
source mirror is false \cite[lines~1483--1492]{Ji2020MIPStar}.  If degrees at
least $q$ are allowed, equality as a function would not by itself prove encoding
membership either, since $x^q$ and $x$ induce the same function on $\mathbb F_q$.

\section{Required repair}

The extraction proof must split the marginal measurement into encoding and
non-encoding outcomes.  On encoding outcomes, the restricted identity above
justifies the substitution.  On the complement, the proof must bound
\[
  \sum_{g\ \mathrm{not\ an\ encoding}}
  \langle\widehat\psi,\widehat S_g^W\widehat\psi\rangle.
\]
The blueprint derives an $O(\delta_S)+md/q$ bound from point consistency and
the combined point measurement.  A non-encoding polynomial differs from every
encoded polynomial, so Schwartz--Zippel bounds the probability that their
evaluations coincide.  This estimate is the open reduction; it must be proved
from the paper hypotheses rather than inserted as an extraction hypothesis.

The printed reindexing at lines~1805--1822 replaces
$g_{\operatorname{Dec}(g_W)}$ by $g_W$, which is the invalid reverse identity
for an arbitrary polynomial outcome.  In the corrected calculation, the
constraint is
$g_h(u)=\operatorname{Dec}(g_W)\mathbin{\cdot}\operatorname{ind}_m(u)-a$.
After setting $h'=h+\operatorname{Dec}(g_W)$, the defining evaluation formula
$g_k(u)=k\mathbin{\cdot}\operatorname{ind}_m(u)$ and characteristic two turn
that constraint into $g_{h'}(u)=a$.  Thus this step uses neither decoder
linearity nor the identity $\operatorname{Dec}(g_h)=h$
\cite[lines~1805--1822]{Ji2020MIPStar}.

\section{Blueprint and formalization status}

Definition~\path{def:qld-full-field-decoder} and
Lemmas~\path{lem:qld-decoder-linearity} and
\path{lem:qld-decoder-evaluation} in
\path{blueprint/src/chapter/ch16_qpbt_extraction.tex} state the full-field
decoder, its linearity and inverse, and the evaluation identity restricted to
encoding outcomes.  Lemma~\path{lem:qld-nonencoding-mass-bound}
separately records the estimate on the remaining outcomes.

The declarations \path{decodeFq_add}, \path{decodeFq_smul}, and
\path{decodeFq_lowDegreeEncoding} in
\path{MIPStarRE/QPBT/Algebra/Decoding.lean} match the three algebraic assertions.
In the same file, \path{IsEncoding} is the equality
$g_{\operatorname{Dec}(g)}=g$ of polynomial representatives, and
\path{decodeFq_dotProduct_indicatorVec} requires an explicit hypothesis
of that predicate.  It therefore matches the corrected evaluation lemma and
does not assert the false unrestricted identity.  Finally,
\path{nonencodingMarginalMass_le} in
\path{MIPStarRE/QPBT/Extraction/MeasurementConsistency.lean} concludes the complementary
mass bound from the projective strategy and global polynomial-pair witness; it
does not add the bound as a witness field or theorem hypothesis.

The decoder calculations are proved.  \path{lowDegreeEncoding_mem_poly},
\path{decodeFq_add}, \path{decodeFq_smul}, \path{decodeFq_lowDegreeEncoding},
and \path{decodeFq_dotProduct_indicatorVec} carry complete proofs: the degree
side-condition on $g_h$, the two linearity laws, the inversion
$\operatorname{Dec}_{\mathbb F_q}(g_h)=h$, and the evaluation identity under the
explicit \path{IsEncoding} hypothesis.  The complementary estimate
\path{nonencodingMarginalMass_le} remains an open proof obligation.
Issue~\#19 owns the Stage~4.2 statements and this discrepancy record;
\ghissue{47} owns the Stage~4.3 decoder calculations, now discharged, and the
non-encoding support estimate.  No unrestricted theorem, residual assumption, or
additional witness field was introduced.

\section{Supplied measurements and the source construction}

Fix the expanded state $\widehat\psi$ of a projective strategy passing the
Pauli basis test with error $0\leq\epsilon\leq1$. Let
$\widehat S^{\mathrm A}$ and $\widehat S^{\mathrm B}$ be projective
polynomial-pair measurements whose evaluated marginals are consistent with
the opposite expanded point measurements at error $\delta_G\geq0$.
Lemma~\texttt{qld-4-7} constructs such measurements with universal constants
$a>1$ and $0<b<1$ at error
$\delta_G=a(md)^a(\epsilon^b+q^{-b}+2^{-bmd})$
\cite[lines~1267--1274]{Ji2020MIPStar}. Assuming the pair is additional proof
data when compared with this source construction.

For a supplied pair, observable self-consistency follows by measuring the
difference polynomial $g-g_h$, comparing its evaluations with the original
point measurements, and applying Schwartz--Zippel. Equal difference
polynomials have equal decoded labels, since
$\operatorname{Dec}(g-g_h)=\operatorname{Dec}(g)-h$.
The argument therefore uses decoder linearity and the valid inverse identity,
without applying the false unrestricted evaluation identity.
Agreement and two squared-distance triangle inequalities, followed by the
fixed-coordinate character conversion, give the observable squared-distance
bound
\[
 48\delta_G+32N\epsilon+4md/q
 \leq (48+32N)(\delta_G+\sqrt\epsilon+md/q),
\]
where $N$ is the parameter-independent number of Pauli test edges.
This estimate is proved in
\path{tildeObs_selfConsistent_ofGlobalPairWitness}, in
\path{MIPStarRE/QPBT/Extraction/Consistency.lean}. It concerns the same uniform
register average and fixed coordinate as the source
\cite[lines~1476--1605]{Ji2020MIPStar}. It does not construct the pair.

The separate source obligation
\path{exists_pulled_apart_consistency}, in
\path{MIPStarRE/QPBT/Extraction/Construction.lean}, requires one constructed
pair satisfying both point estimates and observable self-consistency. Its
universal constants precede the strategy quantifiers. The proof hole records
the source realignment in \ghissue{519}. Discharge requires
\path{exists_globalPairWitness}, including zero error: the current constructor
in \path{MIPStarRE/QPBT/Combining/Apply.lean} is unfinished and assumes strictly
positive error, whereas the source setting includes zero
\cite[lines~197--199]{Ji2020MIPStar}.

Both point estimates remain unfinished here under the names
\path{tildeM_consistent_pointMeas_ofGlobalPairWitness} and
\path{tildeM_consistent_pointMeas'_ofGlobalPairWitness}, in
\path{MIPStarRE/QPBT/Extraction/MeasurementConsistency.lean}. Once established,
they can be applied to the same constructed pair as the observable estimate;
the maximum of the three universal constants supplies a common scale.
The source blueprint node and the separate conditional entry have no
completion marker. A proved supplied-witness estimate does not discharge
the construction premise.

\section{Conclusion}

The decoding identity is false for a general polynomial outcome.  Restricting
it to the encoding image is a necessary local correction.  Recovering the
source's measurement-consistency conclusion then requires the separate
non-encoding-mass estimate and the global measurement construction, which
remain open here. Observable self-consistency is proved for supplied global
measurements; it does not resolve either point estimate or the joint source
construction. The additional witness premise is formal proof debt, not an
assumption of the paper theorem.

\bibliographystyle{alpha}
\bibliography{references}

\end{document}
